\section{Data Integrity dan Non-Repudiation}\label{studi-data-integrity-non-repudiation}

Integrasi \textit{patient-generated health data} (PGHD) ke dalam sistem rekam medis elektronik terdistribusi menimbulkan kebutuhan terhadap mekanisme kriptografi yang dapat menjamin dua properti data, yaitu integritas data dan \textit{non-repudiation}.
Berbeda dengan data klinis yang dicatat oleh tenaga kesehatan, PGHD berasal langsung dari pasien melalui perangkat atau input manual, sehingga kepercayaan terhadap keaslian dan keutuhan data tersebut tidak dapat diasumsikan begitu saja \parencite{Khatiwada_Yang_Lin_Blobel_2024_PGHD_Understanding,Abdolkhani_2019_PGHD_Quality_Challenges}.

\subsection{Data Integrity}\label{studi-data-integrity}

Integritas data (\textit{data integrity}) adalah jaminan bahwa data, konsisten, akurat dan tidak dimodifikasi secara tidak sah.
Integritas merupakan salah satu dari tiga pilar utama yang dikenal sebagai CIA (\textit{confidentiality}, \textit{integrity}, \textit{availability}) \parencite{tripathi2025cia}.

Mekanisme utama untuk menjamin integritas data adalah fungsi hash kriptografis \parencite{tripathi2025cia}.
Standar NIST mendefinisikan fungsi hash kriptografis sebagai fungsi dengan keluaran berukuran tetap yang dirancang memiliki sifat one-way dan \textit{collision resistant}, sehingga sulit secara komputasional untuk mencari masukan yang memetakan ke keluaran tertentu ataupun sepasang masukan berbeda dengan hash yang sama \parencite{NIST_FIPS186-5}.
Dengan demikian, apabila penerima data menghitung ulang hash dari data yang diterima dan hasilnya sesuai dengan hash yang dikirim bersama data, maka data dinyatakan belum mengalami modifikasi.

Meskipun hash kriptografis dapat menjamin integritas, mekanisme ini saja tidak cukup untuk memastikan bahwa data berasal dari sumber yang sah.
Penyerang yang memodifikasi data juga dapat mengganti nilai hash secara bersamaan.
Oleh karena itu, integritas data dalam konteks PGHD harus dikaitkan dengan mekanisme yang sekaligus membuktikan identitas pengirim, yaitu \textit{non-repudiation}.

\subsection{Non-Repudiation}\label{non-repudiation}

\textit{Non-repudiation} adalah layanan kriptografi yang digunakan untuk memberikan jaminan terhadap integritas dan asal-usul data sedemikian rupa sehingga integritas dan asal-usul tersebut dapat diverifikasi dan divalidasi oleh pihak ketiga sebagai berasal dari entitas tertentu yang memiliki kunci privat (yakni penanda tangan) \parencite{NIST_glossary_nonrepudiation,banerjee2025blockchain}.

Dalam konteks PGHD, \textit{non-repudiation} menjadi krusial karena data yang masuk ke dalam sistem RME harus dapat diatribusikan secara pasti kepada pasien yang bersangkutan.
Tanpa mekanisme ini, penyedia layanan kesehatan tidak memiliki bukti kriptografis yang dapat memverifikasi bahwa entri PGHD memang berasal dari pasien yang mengklaimnya \parencite{Khatiwada_Yang_Lin_Blobel_2024_PGHD_Understanding}.

\subsection{Mekanisme Kriptografi untuk Integritas dan Non-Repudiation}\label{sec:mekanisme-kriptografi}

Terdapat beberapa mekanisme kriptografi yang dapat digunakan untuk menjamin integritas data dan \textit{non-repudiation}, dengan tingkatan jaminan dan karakteristik yang berbeda.

\subsubsection{Hash Function}
Fungsi hash kriptografis adalah fungsi satu arah yang menghasilkan \textit{digest} berukuran tetap dari masukan berukuran sembarang \parencite{lei2023analysis,NIST_FIPS186-5}.
Sifat deterministik namun tidak dapat dibalikkan menjadikan nilai sebelumnya sebagai komponen inti untuk verifikasi integritas.
Namun, fungsi hash yang berdiri sendiri tidak memberikan jaminan \textit{non-repudiation} karena tidak mengikat identitas pengirim pada data \parencite{banerjee2025blockchain}.

\subsubsection{Message Authentication Code (MAC)}
\textit{Message Authentication Code} (MAC) adalah nilai autentikasi kecil yang dihasilkan dari kombinasi pesan dan kunci rahasia simetris.
MAC dapat menjamin integritas dan autentisitas data antara dua pihak yang berbagi kunci.
Namun, karena kunci bersifat simetris (dimiliki oleh kedua pihak), MAC tidak dapat memberikan jaminan \textit{non-repudiation}, karena pihak mana pun dari kedua pihak dapat menghasilkan MAC yang valid untuk pesan yang sama, sehingga tidak ada bukti yang dapat membuktikan pihak mana yang memodifikasi data \parencite{NIST_glossary_nonrepudiation}.

\subsubsection{Digital Signature}
\textit{Digital signature} adalah mekanisme untuk memverifikasi autentisitas asal data, integritas data, dan \textit{non-repudiation} penanda tangan \parencite{NIST_glossary_digitalsig}.
Berbeda dengan MAC, \textit{digital signature} memanfaatkan kriptografi asimetris, di mana pengirim menandatangani data menggunakan kunci privat yang hanya dimilikinya, sedangkan verifikasi dilakukan oleh siapa pun menggunakan kunci publik yang bersesuaian.

\textit{Digital signature} digunakan untuk mendeteksi modifikasi tidak sah terhadap data dan untuk mengautentikasi identitas penanda tangan. % chktex 8
Di samping itu, penerima data yang telah ditandatangani dapat menggunakan \textit{digital signature} sebagai bukti kepada pihak ketiga bahwa tanda tangan tersebut memang dibuat oleh penanda tangan yang diklaim.
Sifat inilah yang dikenal sebagai \textit{non-repudiation}, karena penanda tangan tidak dapat dengan mudah menyangkal tanda tangan tersebut di kemudian hari \parencite{NIST_FIPS186-5}.

% NIST FIPS 186-5 menspesifikasikan tiga algoritma utama untuk \textit{digital signature}, yaitu RSA, \textit{Elliptic Curve Digital Signature Algorithm} (ECDSA), dan \textit{Edwards Curve Digital Signature Algorithm} (EdDSA) \parencite{NIST_FIPS186-5}. % chktex 8
% RSA merupakan skema tanda tangan berbasis masalah faktorisasi bilangan bulat, di mana kunci publik terdiri atas modulus $n = pq$ dan eksponen publik $e$, sedangkan tanda tangan dihasilkan menggunakan eksponen privat $d$ \parencite{NIST_FIPS186-5}.
% ECDSA adalah analog DSA pada kriptografi kurva eliptik, yang keamanannya bergantung pada kesulitan \textit{elliptic-curve discrete logarithm problem} (ECDLP) dan menggunakan pasangan kunci pada kurva yang direkomendasikan dalam SP~800-186 \parencite{NIST_FIPS186-5}.  % chktex 8
% EdDSA, sebagaimana distandarkan dengan kurva edwards25519 dan edwards448, merupakan skema tanda tangan deterministik berbasis kurva Edwards yang dirancang untuk efisiensi tinggi dan keamanan sebanding dengan panjang kunci yang lebih kecil \parencite{NIST_FIPS186-5}.

% \subsubsection{Public Key Infrastructure (PKI)}

% Public Key Infrastructure (PKI) adalah himpunan kebijakan, peran, perangkat lunak, dan prosedur yang digunakan untuk mengelola siklus hidup pasangan kunci publik--privat dan sertifikat digital, termasuk penerbitan, distribusi, penggunaan, serta pencabutan sertifikat \parencite{NIST_PKI}. % chktex 8
% Dalam PKI, Certificate Authority (CA) dipercaya menerbitkan sertifikat digital yang mengikat identitas entitas (misalnya pengguna, perangkat, atau layanan) dengan kunci publik yang bersesuaian, sehingga menyediakan CA berperan sebagai regulator layanan autentikasi, integritas, dan \textit{non-repudiation} \parencite{NIST_PKI,NIST_SP80057}.

% \subsubsection{Certificateless Signature}
% \textit{Certificateless signature} adalah varian skema tanda tangan digital yang menghilangkan kebutuhan akan manajemen sertifikat PKI (\textit{Public Key Infrastructure}) sekaligus mengatasi masalah \textit{key escrow} yang melekat pada skema berbasis identitas \parencite{kakkar2021certificateless}.
% Dalam skema ini, kunci publik pengguna tidak memerlukan sertifikat dari otoritas terpercaya untuk membuktikan keasliannya, sehingga dapat dioperasikan tanpa CA (\textit{Certificate Authority}) terpusat.
